\documentclass[11pt]{article}
\usepackage[margin=1in]{geometry}
\usepackage{amsmath,amssymb,amsthm,booktabs,hyperref}
\newtheorem{proposition}{Proposition}
\newtheorem{lemma}{Lemma}
\title{Torsion from Wilson-line lattice shifts in heterotic orbifolds:\\[2pt]
\large a tested phase formula, its selection rule, and a Scherk--Schwarz window}
\author{W(3,3) programme --- Claude track (Passes 12265--12333)}
\date{October 2026}
\begin{document}
\maketitle

\begin{abstract}
Shifting a Wilson line of a heterotic orbifold by a lattice vector $\lambda$ changes no consistency condition but can change the spectrum. We
write the change of every projection phase as an alternating bicharacter $B(g,h)$ on commuting space-group elements and test it against
the orbifolder engines: $B=0$ on all commuting pairs implies an unchanged spectrum ($232/232$), every change has $B\neq0$ ($61/61$), and a
pre-registered prediction for fresh $T^6/\mathbb Z_4$ models is correct in $124/124$ cases. Brother models and generalized discrete
torsion are known (Pl\"oger, Ramos-S\'anchez, Ratz, Vaudrevange); what is new here is the explicit, tested formula, its use on the
$W(3,3)$ orbifold routes, and a radius analysis of a Scherk--Schwarz brother with three chiral generations.
\end{abstract}

\section{The phase formula}
Let $\lambda\in\Gamma_{E_8\times E_8}$ be added to the Wilson lines in a set $S$ of lattice directions, and write $\sigma_g=\sum_{x\in S}m_x(g)$
for the winding labels of a space-group element $g$ along $S$. In a sector $h$ the left-moving momenta form the coset
$\Gamma+V_h$, which the shift does not change.
\begin{lemma}
The projection phase of $g$ on sector $h$ changes by
\[
B(g,h)=\tfrac12\bigl(\sigma_g\,V_h\cdot\lambda-\sigma_h\,V_g\cdot\lambda\bigr)\pmod 1,
\]
an alternating bicharacter on commuting pairs ($\lambda^2\in2\mathbb Z$ makes the quadratic term integral).
\end{lemma}
\emph{Tests.} On $320$ relabellings ($T^6/\mathbb Z_3$ and $T^6/\mathbb Z_2\times\mathbb Z_2$), $B$ evaluated on every commuting pair
predicts every spectrum change; $B\neq0$ is necessary, not sufficient (three cases with $B\neq0$ leave every massless state unchanged).
A first version that dropped fixed-point labels mispredicted $16/100$ $T^6/\mathbb Z_4$ brothers; with them, predictions written and pushed
before the engine run were correct in $124/124$ cases.

\section{Consequences}
\begin{proposition}
On $T^6/\mathbb Z_3$ a lattice shift preserving the Pl\"oger conditions induces no torsion: every twisted centralizer is generated by
the sector element and the central fermion-parity element, on which $B$ vanishes by antisymmetry and by the conditions themselves.
\end{proposition}
This explains why the non-supersymmetric orbifolder, which skips two of the conditions (a verbose-only failure flag), still produced
correct $T^6/\mathbb Z_3$ spectra: all $497$ affected corpus inputs have a lattice-shift repair with an identical spectrum ($983/983$).
Torsion requires commuting elements whose effective orders share a factor, and it also enters through fixed-point labels.
\begin{itemize}
\item Supersymmetric $\mathbb Z_6$-II Standard Models: $1471/1471$ brothers identical.
\item $\mathbb Z_2\times\mathbb Z_6$-I: $12$ further three-generation supersymmetric Standard Models. All fail later tests: renormalisable
$QLd^c$, or textures.
\item Non-supersymmetric $\mathbb Z_6$-II: the fermion-parity element and an order-2 translation give a Scherk--Schwarz bit.
\end{itemize}

\section{A Scherk--Schwarz brother with three chiral generations}
Among all $1399$ non-supersymmetric $\mathbb Z_6$-II models of the route, exactly one lineage ($15$ variants of one parent) has a
Scherk--Schwarz brother with three chiral left-right generations and no chiral fractional charge. Its brother has no tachyon at zero
momentum and winding (engine solver). The torsion shifts the momentum of the fermion-parity-twisted states by $1/2$ for every winding (a
lattice Wilson line is not a freely acting translation). Hence:
\begin{itemize}
\item it is tachyonic for $R_6>\sqrt{3\alpha'/8}$;
\item exact $E_8\times E_8$ level minima ($-1/3$ against $-2/3$ needed, $0$ against $-1/2$) close every small-radius channel;
\item a winding channel along a second circle requires $R_5\gtrsim\sqrt{2\alpha'}$.
\end{itemize}
\emph{Computed:} zero-winding spectra, right-moving tachyonic vacua, the momentum shift, the level minima. \emph{Argued:} the circle mass
formula; the projections of winding states are not computed.

\paragraph{Obstructions.} The hypercharge normalisation is $k_Y=8/3$. One-loop running with the model's full light spectrum above any
common threshold between $10^3$ and $10^{16}$~GeV needs $k_Y\in[0.21,1.28]$ at a perturbative $\alpha_2=\alpha_3$ crossing. The one-loop
potential in the string-scale window needs the states the momentum shift moves to half-integer momentum, which the engine does not
produce. The candidate is therefore not a Standard Model.

\paragraph{Record of corrections.} The radius statement was revised twice before reaching the form above:
\begin{itemize}
\item winding instead of momentum;
\item an odd-winding flip that holds only for free actions.
\end{itemize}
A suspicion that the engine's ``excited tachyons'' were massless was withdrawn after its solver was run verbosely. Each correction is
recorded in the producers of Passes 12299--12327.

\paragraph{Reproducibility.} Producers \texttt{w33\_pass12265\ldots} through \texttt{w33\_pass12328\ldots} in \texttt{analysis/}; the
pre-registration is commit \texttt{e570f4b1f}.
\end{document}
